\documentclass[11pt]{article}

\usepackage[a4paper,margin=1in]{geometry}
\usepackage{amsmath,amssymb}
\usepackage[hidelinks]{hyperref}
\usepackage{xcolor}
\usepackage{enumitem}
\usepackage{listings}

\lstset{
  basicstyle=\ttfamily\small,
  breaklines=true,
  columns=fullflexible,
  frame=single,
  backgroundcolor=\color{gray!6}
}

\title{Neue Anmerkungen zur verbesserten Fassung\\
\large Review-Notizen zum arXiv-Manuskript}
\author{Kommentare von ChatGPT}
\date{\today}

\begin{document}
\maketitle

\section*{Gesamturteil}

Die neue Fassung ist erheblich besser und wirkt deutlich publikationsnäher. Die wichtigsten früheren Schwachstellen sind behoben:

\begin{itemize}[leftmargin=2em]
    \item Der Hauptsatz wird jetzt sauber über die Proposition zum generischen Fall und den Satz zum degenerierten Fall bewiesen.
    \item Die Integrabilität von \(y\) ist ergänzt.
    \item Die Minimalität im generischen Fall ist als eigene Proposition isoliert.
    \item Das Lemma über den trivialen Stabilisator eines zyklischen Intervalls ist explizit formuliert.
    \item Die Set/Multiset-Teilbarkeitsaussage ist sauber in das Korollar aufgenommen.
    \item Die Korollar-Sektion ist nun klar als Multiset-Sektion markiert.
    \item Die Lean-Zeilenzahl ist konsistenter gemacht.
    \item Das Beispiel \(N>D\), \(D\nmid N\) ist viel anschaulicher geworden.
\end{itemize}

Mein Urteil: Die Fassung ist in der Beweisarchitektur deutlich näher an einer arXiv-Version. Es bleiben aber einige Punkte, die vor einer Einreichung verbessert werden sollten.

\section{Degenerierter Fall: Minimalität noch explizit beweisen}

Im Satz zum degenerierten Fall wird gezeigt, dass die Gap-Sequenz die Form
\[
    \underbrace{0,\ldots,0}_{N/D-1},1,
    \underbrace{0,\ldots,0}_{N/D-1},1,\ldots
\]
hat. Damit ist klar, dass \(N/D\) eine Periode ist. Der Satz behauptet aber die \emph{minimale} Periode. Diese Minimalität sollte ausdrücklich begründet werden.

Empfohlene Ergänzung direkt nach der dargestellten Gap-Sequenz:

\begin{lstlisting}[language={[LaTeX]TeX}]
If \(N/D=1\), this period is clearly minimal. If \(N/D>1\), the value
\(1\) occurs exactly once in each block of length \(N/D\), while all other
entries in the block are \(0\). Hence no smaller positive period can
preserve the positions of the entries equal to \(1\). Therefore the
minimal period is \(N/D\).
\end{lstlisting}

Das ist der wichtigste noch offene mathematische Präzisierungspunkt.

\section{Minimalitätsbeweis: Blocksumme präziser formulieren}

Im generischen Minimalitätsbeweis wird definiert:
\[
    \sigma := \delta^{(0)}+\delta^{(1)}+\cdots+\delta^{(\lambda'-1)}.
\]
Danach folgt unmittelbar:
\[
    \tilde p_s^{(i+\lambda')} = \tilde p_s^{(i)} + \sigma.
\]
Das ist richtig, aber ein Zwischensatz würde die Argumentation wasserdicht machen.

Empfohlene Ergänzung:

\begin{lstlisting}[language={[LaTeX]TeX}]
By \(\lambda'\)-periodicity, every block of \(\lambda'\) consecutive gaps
has the same sum \(\sigma\); that is, for every \(i\in\mathbb Z\),
\[
    \delta^{(i)}+\delta^{(i+1)}+\cdots+\delta^{(i+\lambda'-1)}=\sigma.
\]
Hence
\[
    \tilde p_s^{(i+\lambda')}-\tilde p_s^{(i)}=\sigma.
\]
\end{lstlisting}

\section{Translation invariance des Multisets präzisieren}

Die Aussage, dass aus
\[
    \tilde p_s^{(i+\lambda')} = \tilde p_s^{(i)}+\sigma
\]
die Translationsinvarianz des gesamten Multisets folgt, ist korrekt. Sie sollte aber einen Satz mehr bekommen, damit klar ist, dass Multiplikitäten erhalten bleiben.

Empfohlene Ergänzung:

\begin{lstlisting}[language={[LaTeX]TeX}]
Since \(i\mapsto i+\lambda'\) is a bijection of \(\mathbb Z\), this
translation maps the enumerated multiset bijectively onto itself and
preserves multiplicities.
\end{lstlisting}

\section{Formulierung zur schweren Menge leicht verbessern}

Die aktuelle Formulierung, dass die ``final \(s\) elements'' des Blocks den Extra-Hit liefern, ist inhaltlich richtig, aber etwas missverständlich. Nach \(q\) vollständigen Umläufen sind die verbleibenden \(s\) Elemente entscheidend; man sollte diese Struktur direkt so formulieren.

Empfohlene Ersatzformulierung:

\begin{lstlisting}[language={[LaTeX]TeX}]
because after \(q\) complete cycles the remaining \(s\) elements contribute
one additional hit on \(s\) consecutive classes in the \(m\)-stepping order.
\end{lstlisting}

\section{Lean-Sektion: Umfang noch früher präzisieren}

Die spätere Scope-Erklärung ist gut. Noch besser wäre es, bereits am Anfang der Lean-Sektion deutlich zu sagen, dass die geometrische Konstruktion nicht vollständig von Grund auf formalisiert wurde, sondern nach Reduktion auf das kombinatorische Restklassenmodell.

Empfohlener Satz am Anfang der Lean-Sektion:

\begin{lstlisting}[language={[LaTeX]TeX}]
The formalisation verifies the residue-class and period-minimality
arguments after the geometric construction has been reduced to the
combinatorial residue model.
\end{lstlisting}

Im Abstract könnte der Satz

\begin{lstlisting}[language={[LaTeX]TeX}]
The algebraic core of both proofs has been formalised in Lean~4.
\end{lstlisting}

leicht präzisiert werden zu:

\begin{lstlisting}[language={[LaTeX]TeX}]
The algebraic and combinatorial core of both proofs has been formalised in Lean~4.
\end{lstlisting}

\section{Acknowledgements entschärfen}

Der aktuelle Hinweis auf KI-Unterstützung ist transparent, aber die Formulierung ``proof construction'' könnte unnötig Misstrauen erzeugen. Professioneller und nüchterner wäre:

\begin{lstlisting}[language={[LaTeX]TeX}]
The author used ChatGPT, Gemini, and Claude as assistive tools during
mathematical exploration, drafting, code generation, and proofreading.
The author has independently verified the mathematical content and accepts
full responsibility for the results and any errors.
\end{lstlisting}

\section{Abschließende Einschätzung}

Nach den bisherigen Verbesserungen bleibt als wichtigste mathematische Ergänzung nur noch die explizite Minimalität im degenerierten Fall. Danach ist das Manuskript nach meiner Einschätzung arXiv-tauglich als elementarer, klarer Beitrag über rationale Cut-and-Project-Strip-Projektionen, vorausgesetzt Repository, Lean-Code und empirische Daten sind tatsächlich konsistent mit den Angaben.

Vor Einreichung sollte das Dokument außerdem einmal lokal kompiliert werden, insbesondere wegen Label-Umbenennungen wie \texttt{subsec:gapseq} und wegen der Referenzen auf die Lean-Tabelle.

\end{document}
